\documentclass[12pt]{article}

% Layout.
\usepackage[top=1.2in, bottom=0.9in, left=1in, right=1in, headheight=1in, headsep=6pt]{geometry}

% Fonts.
\usepackage{mathptmx}
\usepackage[scaled=1.0]{helvet}
\renewcommand{\emph}[1]{\textsf{\textbf{#1}}}

% Misc packages.
\usepackage{amsmath,amssymb,latexsym}
\usepackage{graphicx,hyperref}
\usepackage{array}
\usepackage{xcolor}
\usepackage{multicol}
\usepackage{tabularx,colortbl}
\usepackage{enumitem}
\usepackage[nodayofweek]{datetime}

\hypersetup{
    colorlinks=true,
    linkcolor=blue,
    filecolor=magenta,      
    urlcolor=blue,
    pdftitle={Syllabus for MATH F252X section 901 Fall 2026 (Meek)},
    }

\def\mailto#1{\href{mailto:#1}{#1}}

% Paragraph spacing
\parindent 0pt
\parskip 6pt plus 1pt
\def\tableindent{\hskip 0.5 in}
\def\ts{\hskip 1.5 em}

\usepackage{fancyhdr}
\pagestyle{fancy} 
\lhead{\large\sf\textbf{MATH F252 Calculus II}}
\chead{\large\sf\textbf{Syllabus}}
\rhead{\large\sf\textbf{Fall 2026}}

\newcommand{\localhead}[1]{\par\smallskip\textbf{#1} \smallskip\nobreak\\}%
\def\heading#1{\localhead{\large\emph{#1}}}
\def\subheading#1{\localhead{\emph{#1}}}

\newenvironment{clist}%
{\bgroup\parskip 0pt\begin{list}{$\bullet$}{\partopsep 4pt\topsep 0pt\itemsep -2pt}}%
{\end{list}\egroup}%

\begin{document}

\strut\par\vskip-12pt
\heading{Essential Information}

\vskip -12pt
\strut\hbox to \hsize{\tableindent\vtop{\halign{#\hfill\ts&#\hfil\cr
{\emph{Instructor}} & Kevin Meek \quad \href{mailto:krmeek2@alaska.edu}{\texttt{krmeek2\@@alaska.edu}} \cr
\strut & \cr
{\emph{Public webpage}}&\href{https://uaf-math.github.io/calc2/}{\texttt{https://uaf-math.github.io/calc2}}\cr
\strut & \cr
{\emph{Canvas webpage}} & \href{https://canvas.alaska.edu/courses/31682}{\texttt{canvas.alaska.edu/courses/31682}} \cr
\strut & \cr
\emph{Prerequisite} & MATH F251X Calculus I; or placement.\cr
\strut & \cr
{\emph{Required text}} &\textit{OpenStax Calculus Volume 2} by G. Strang \& E. Herman,\cr
&\href{https://openstax.org/details/books/calculus-volume-2}{\texttt{openstax.org/details/books/calculus-volume-2}} \cr
&
  optional paperback print copy:\quad \texttt{ISBN-13 978-1-50669-807-6}\cr
  }
\hfil}}


\heading{Description and Course Goals}
Calculus is useful in all of the science and engineering disciplines,
so it is part of the UAF core curriculum.  The two principal tools
from Calculus I (MATH 251), differentiation and integration, are
central here too.  However, this course extends your understanding of
integration especially.  We will also study sequences and series,
Taylor series, and parametric and polar curves.

At the completion of the course, students will:
\begin{clist}
\item possess mature skills in computing integrals in one variable,
\item be able to apply integrals to common problems (areas, volumes, work, \dots),
\item understand the convergence of sequences and series,
\item understand and be able to apply Taylor series and polynomials, and
\item be able to use parametric and polar equations for curves.
\end{clist}
Upon completion, students will have the mathematical foundation for
success in Calculus III and other courses requiring mathematics
background, including many junior/senior-level science and engineering
courses.

%FIXME FROM HERE
\heading{Class Time}
There are 4.5 hours of class meetings with your instructor every week:
\begin{itemize}
\item MWF 11:45 -- 12:45 Gruening 208
\item Th 11:30 -- 13:00 Gruening 208
\end{itemize}

The 1.5 hour time on Thursday will be used to discuss new topics (30
minutes) and for a weekly Quiz (30 minutes).  The final 30 minutes, at
most, will be spent going over the Quiz.  We will go over the quiz
immediately so that students can leave knowing what Quiz material has
been mastered, how to work each problem correctly, and what topics
need additional work.


\clearpage\newpage

\strut

\vspace{-12pt}

\heading{Schedule and Online Materials}
The course website contains a
\href{https://uaf-math.github.io/calc2/assets/general/S25/meek/schedule.pdf}{Schedule}
schedule
listing the textbook sections to be covered each class, the dates each
Homework is due, plus the dates for Quizzes and Exams. You should
consult this schedule frequently.  I will announce Schedule
adjustments in class.

Most course materials (Syllabus, old Quiz and Exam solutions, study materials, etc.) will be posted on the \href{https://uaf-math.github.io/calc2/}{course webpage}.  Some course materials (grades, Homework solutions, announcements, etc.) will be available on the \href{https://canvas.alaska.edu/courses/31682}{Canvas site}.  Each website links to the other.


\heading{Office Hours and Communication}

I will use \href{https://canvas.alaska.edu/courses/31685}{Canvas} to
send announcements, and, if I need to contact you outside of the
lecture then I'll try to email via Canvas.  (\textsl{Please set the
  email address in Canvas to one that you check regularly!})

My office hours are Mondays, Wednesdays,
and Fridays 9:15-10:15am in Chapman 301C and Thursdays 10:15-11:15 in
the
Math and Stats Lab. If you are unable to make my
office hours, I am happy to meet by appointment at another
time.

For e-mail communication, please always include your course and
section number in your e-mail. The easiest way to do this is to
message me through Canvas, as this will automatically send me an
e-mail with the relevant information included.

Please let me know by email if you will miss a class for an excusable reason.

\heading{Evaluation and Grades}
Grades are determined as follows.  (Each component of the grade is discussed below.)
 
\begin{multicols}{2}
\begin{tabular}{|c|c|}
\hline
Participation/attendance & 5\%\\
\hline
Homework & 15\% \\
\hline
Quizzes & 20\% \\
\hline
Midterm Exam 1 & 20\% \\
\hline
Midterm Exam 2 & 20\%  \\
\hline
Final Exam & 25\% \\
\hline
total & 100\% \, \\
\hline
\end{tabular}

%\vskip 6pt

\begin{tabular}{llll}
A  & 93--100\%& C  & 68--75\%  \\
A- & 90--92.99\% & C- & not given \\
B+ & 87--89.99\% & D+ & 65--67.99\%  \\
B  & 82--86.99\% & D  & 60--64.99\%  \\
B- & 79--81.99\% & D- & 57--59.99\%  \\
C+ & 76--78.99\% & F  & $\le$ 56.99\%
\end{tabular}
\end{multicols}

The grade ranges above are a guarantee and a lower bound. I reserve
the right to increase your grade above these ranges based on the
actual difficulty of the work and/or on average class performance.
Any such increases will preserve grade ordering by weighted total
score.


\heading{Participation}
Attendance and participation are important parts of mastering the
material, and a strong predictor of overall course performance.
Attendance, class participation, and participation in groupwork will
count towards this portion of the grade.

\newpage

\heading{Homework}
Homework assignments use a selection of problems from the textbook.
Please write your Homework solutions on paper, or electronically on a
tablet, and turn it in as a PDF via Gradescope.  Gradescope is
accessed via the \href{https://canvas.alaska.edu/courses/31682}{Canvas
page}.  Help with scanning your Homework can be found on the
\href{https://uaf-math.github.io/calc2/techHelp.html}{Tech Help}
webpage.  Assignments are due by 11:59pm on the date assigned; see the
online 
schedule
% \href{https://uaf-math.github.io/calc2/assets/general/S24/Bueler/schedule.pdf}{Schedule}
for due dates (currently visible on Canvas, coming soon on the public webpage).  The list of problems is at the
\href{https://uaf-math.github.io/calc2/homework.html}{Homework} tab
on the public webpage.

\emph{Complete worked solutions to all Homework problems are provided
in advance on the Canvas site.}  Therefore Homework will be graded
based on \emph{effort} \emph{completion}, \emph{organization}, and
\emph{neatness}.  Illegible homework assignments will receive a
zero. All students should
earn 100\% on homework.  Of course, it is possible, to defeat the
purpose of the homework by copying the solutions.  This is a bad idea,
and because Homework is only 10\% of your grade it is not worthwhile.
The Homework exists so you can learn by doing! Gradescope is
configured to accept homework up to two days after the due date. After
that, late homework will not be accepted under any circumstances. It
is your responsibility to seek help and complete homework
assignments on time.


\heading{Quizzes}
A weekly Quiz will be given on Thursdays.  It will be in the middle
third of the 1.5 hour class.  (Thursday classes are 9:45am --
11:15am.)  A Quiz will cover material which appeared on recent
Homeworks, since the previous Quiz.  See the
\href{https://uaf-math.github.io/calc2/quizzes.html}{Quizzes tab on
the public site} for coverage.

Quizzes are given under Exam conditions: books, notes, and calculators
are not allowed.  Performance on Quizzes is your best indicator of how
well you are learning the course material, and it is a much better
predictor of Exam performance than is your Homework score.

However, students will be given the immediate opportunity to grade and
correct their Quizzes in the last third of the Thursday class.  You
can earn-back up to half the missed points by doing so accurately.

% Always contact me if you will miss a Quiz!  There are enough Quizzes
% during the semester so that make-up Quizzes will normally not be
% created.  Instead I might ask a coach to proctor your Quiz if you are on
% the road, for example, or I will mark it as excused.

\heading{Midterm and Final Exams}
There are two Midterm Exams, to be held on the dates in the Schedule on the course website:
\begin{itemize}
\item \emph{Midterm Exam 1 on Thursday, 1 October}
\item \emph{Midterm Exam 2 on Thursday, 12 November}
\end{itemize}
Midterms are given during the class time.

The cumulative \emph{Final Exam} will be held at the day/time listed
as the common math time in the online schedule: \emph{10:15-12:15
Tuesday, 8 December}. Department policy does not allow me to give an
early Final Exam.

All Exams will be closed book, closed notes, and no calculator.

\newpage

\heading{Make-ups}
Make-up exams and quizzes may be provided at the instructor's
discretion only if
the student shows documentation supporting an excused absence. Excused
absences include absences due to illness, travel on University or
career-related business, and cultural or religious
observations. Unless it is an unforseen circumstance, arrangements
must be made in advance; if unforseen, the instructor must be notified
by the following day.

\heading{Time Commitment}
As with most university courses, you should expect to spend at least
2-3 hours
{\bf outside} of class each week for every credit hour. So for this
course, you should plan to spend a minimum of 8-12 hours actually
working homework and examples from the text, asking questions, and
preparing for quizzes and exams. This does not include the time spent
finding your textbook and homework assignment, getting coffee, etc!
Ideally, this time should be spred throughout the week over a minimum
of 5 days.


\heading{Tutoring and Resources}
\vskip -30pt\strut
\begin{clist}
    \item The Math and Stat Lab is now located in the brand new
      student success center on the 6th floor of the Rassmussen Building and
      offers drop-in tutoring.

	See 

	\href{http://www.uaf.edu/dms/mathlab/}{\texttt{www.uaf.edu/dms/mathlab/}} for schedules and availability.
	\item Free
one-on-one (or small group) tutoring is available in 
Chapman Building Room 201. You must schedule an
appointment; see \href{http://www.uaf.edu/dms/mathlab/}{\texttt{www.uaf.edu/dms/mathlab/}}.
	\item Student Support Services (\href{https://uaf.edu/sss/}{\texttt{uaf.edu/sss/}}) offers free tutoring in many subjects to students who qualify for their program.
	\item ASUAF (\href{https://uaf.edu/asuaf/}{\texttt{uaf.edu/asuaf/}}) offers private tutoring for a small fee, based on student income.
\end{clist}



\strut

\vspace{-12pt}

\heading{AI usage}
During proctored and on-paper Quizzes and Exams you will not have access to electronice tools of any type, nor books or notes.  These assessments represent 85\% of your grade.

However, feel free to use a calculator on Homework.  It is also reasonable to explore new AI tools like ChatGPT, but how you guide your prompts and verify the results is what matters.  Merely doing cut and paste without understanding will have no benefit.  Noting that Quizzes and Exams represent the vast majority of your grade, your own thinking, as you do the homework, has the greatest benefits.

\newpage

\heading{Rules and Policies}
\vskip -20pt

\subheading{Incomplete Grade} 
Incomplete (I) will only be given in
  DMS courses in cases where
  the student has completed the majority (normally all but the last
  three weeks) of a course with a grade of C or better, but for
  personal reasons beyond his/her control has been unable to complete
  the course during the regular term. Negligence or indifference are
  not acceptable reasons for the granting of an incomplete. 

\subheading{Late Withdrawals} 
A withdrawal after the deadline from a DMS course will
  normally be granted only in cases where the student is performing
  satisfactorily (i.e., C or better) in a course, but has exceptional
  reasons, beyond his/her control, for being unable to complete the
  course. These exceptional reasons should be detailed in writing to
  the instructor, department head and dean.

\subheading{No Early Final Examinations}
Final examinations for DMS
  courses shall not be held earlier than the date and time published
  in the official term schedule. Normally, a student will not be
  allowed to take a final exam early. Exceptions can be made by
  individual instructors, but should only be allowed in exceptional
  circumstances and in a manner which doesn't endanger the security of
  the exam.

\subheading{Academic Dishonesty}
Academic dishonesty, including cheating and plagiarism, will not
be tolerated.  It is a violation of the Student Code of Conduct
and will be punished according to UAF procedures.

%\begin{center} \textsc{Syllabus Addendum} \end{center}

\subheading{Student protections and services statement}
Every qualified student is welcome in my classroom.  As needed, I am happy to work with you, Disability Services, Veterans' Services, Rural Student Services, etc.~to find reasonable accommodations.  Students at this University are protected against sexual harassment and discrimination (Title IX), and minors have additional protections.  As required, if I notice or am informed of certain types of misconduct, then I am required to report it to the appropriate authorities.  For more information on your rights as a student and the resources available to you to resolve problems, please go the following site: \href{https://www.uaf.edu/handbook/}{\texttt{www.uaf.edu/handbook/}}.

\clearpage\newpage

\strut

\vspace{-12pt}


\subheading{General education statement}
This course is listed as a General Education Math Course.  As such
this course is expected to \textsl{contribute to meeting the
following} four general learning outcomes:

\begin{enumerate}
\item Build knowledge of human institutions, sociocultural processes,
and the physical and natural works through the study of mathematics.
Competence will be demonstrated for the foundational information in
each subject area, its context and significance, and the methods used
in advancing each.

\item Develop intellectual and practical skills across the curriculum,
including inquiry and analysis, critical and creative thinking,
problem solving, written and oral communication, information literacy,
technological competence, and collaborative learning. Proficiency will
be demonstrated across the curriculum through critical analysis of
proffered information, well-reasoned solutions to problems or
inferences drawn from evidence, effective written and oral
communication, and satisfactory outcomes of group projects.

\item Acquire tools for effective civic engagement in local through
global contexts, including ethical reasoning, intercultural
competence, and knowledge of Alaska and Alaska issues.  Facility will
be demonstrated through analyses of issues including dimensions of
ethics, human and cultural diversity, conflicts and interdependencies,
globalization, and sustainability.

\item Integrate and apply learning, including synthesis and advanced
accomplishment across general and specialized studies, adapting them
to new settings, questions and responsibilities, and forming a
foundation for lifelong learning. Preparation will be demonstrated
though production of a a creative or scholarly product that requires
broad knowledge, appropriate technical proficiency, information
collection, synthesis, interpretation, presentation and reflection.
\end{enumerate}

\begin{center}
\textbf{\large{Official UAF Syllabus Addendum}}
\end{center}


\noindent{\bf Student protections statement:} The university respects
and upholds the principles of due process and a fair and equitable
process as specified in the Board of Regents' Policy 09.02 Student
Rights and Responsibilities. For more information regarding the rights
and responsibilities of students, refer to the Office of Rights,
Compliance and Accountability website. You are encouraged to read the
Board of Regents' policy carefully to fully understand your
responsibilities to our community.

We strive to create a safe and respectful environment for all members
of our community. If you have questions about expectations of you as a
student or believe your rights are being violated, we encourage you to
reach out to the Office of Rights, Compliance and Accountability for
help. UAF reserves the right to suspend, expel or take other necessary
and appropriate action in cases where a student is unable or unwilling
to uphold community standards and campus safety.

For more information on your rights as a student and the resources
available to you to resolve problems, please go to the following
site:\\
\url{https://catalog.uaf.edu/academics-regulations/students-rights-responsibilities/}.

\noindent{\bf Disability services statement:} UAF is committed to
creating a learning environment that meets the needs of its diverse
student body.  If you have a disability, or think you might have a
disability, please contact Disability Services (DS) to request a
reasonable accommodation.  Accommodations are not retroactive, so
students should request accommodations as early as possible, since
they might take time to implement.  Students should notify DS at any
time during the semester if adjustments to their communicated
accommodation plan are needed, and must request their accommodations
each semester.  More information is available online at
\url{https://www.uaf.edu/disability}, and students can reach DS at
\mailto{uaf-disability-services@alaska.edu} or 907-474-5655.

\noindent{\bf ASUAF Advocacy \& Engagement:} The Associated Students
of the University of Alaska Fairbanks (ASUAF), UAF's student
government, works to represent student voices and improve the student
experience. Students are encouraged to get involved in shaping
policies, sharing feedback, and helping address issues that impact
campus life. ASUAF also offers support in navigating university
processes and connecting with the right resources when challenges
arise. To get involved or seek support, visit the ASUAF office or
email \mailto{asuaf.office@alaska.edu}.

\noindent{\bf Student Academic Support:}
\begin{itemize}
\setlength\itemsep{0em}
        \item Communication Center (907-474-7007, \mailto{uaf-commcenter@alaska.edu}, Student Success Center, 6th Floor Room 677 Rasmuson Library)
        \item Writing Center (907-474-5314, \mailto{uaf-writing-center@alaska.edu}, Student Success Center, 6th Floor Room 677 Rasmuson Library)
\item UAF Math Services (907-474-7332, \mailto{uaf-traccloud@alaska.edu})


\begin{itemize}
\item Drop-in tutoring, Student Success Center, 6th Floor Room 672 Rasmuson Library

\item 1:1 tutoring (by appointment only), 6th Floor Room 677 Rasmuson Library

\item Online tutoring (by appointment only)
available at the Student Success Center, schedule at \url{https://www.uaf.edu/dms/mathlab/}, 
\end{itemize}

\item Developmental Math Lab, Gruening 406
\item The Debbie Moses Learning Center at CTC (907-455-2860, 604 Barnette St, Room 102,\\ \url{https://www.ctc.uaf.edu/student-services/student-success-center/})
\end{itemize}

 For more information and resources, please see the Academic Advising Resource List (\url{https://www.uaf.edu/advising/students/index.php})

 \newpage

\noindent{\bf Student Resources:}
\begin{itemize}
\setlength\itemsep{0em}
\item Disability Services (907-474-5655, \mailto{uaf-disability-services@alaska.edu}, 110 Eielson Building)
\item Student Health \& Counseling [free counseling sessions available] (907-474-7043, \url{https://www.uaf.edu/chc/appointments.php}, Whitaker Building, Room 206, Health, Safety \& Security Bldg --- same building as Fire and Police)
\item Office of Rights, Compliance and Accountability (907-474-7300, \mailto{uaf-orca@alaska.edu}, 3rd Floor, Constitution Hall)
\item Associated Students of the University of Alaska Fairbanks (ASUAF) or ASUAF Student Government (907-474-7355, \mailto{asuaf.office@alaska.edu}{asuaf.office@alaska.edu}, Wood Center 119)
\end{itemize}

\noindent{\bf Nondiscrimination statement:}
Nondiscrimination statement: The University of Alaska is an equal opportunity/equal access employer, educational institution and provider. The University of Alaska does not discriminate on the basis of race, religion, color, national origin, citizenship, age, sex, physical or mental disability, status as a protected veteran, marital status, changes in marital status, pregnancy, childbirth or related medical conditions, parenthood, sexual orientation, gender identity, political affiliation or belief, genetic information, or other legally protected status. The University's commitment to nondiscrimination, including against sex discrimination, applies to students, employees, and applicants for admission and employment. Contact information, applicable laws, and complaint procedures are included on UA's statement of nondiscrimination available at \url{www.alaska.edu/nondiscrimination}.


For more information, contact: 
\begin{tabular}{l}
UAF Office of Rights, Compliance and Accountability\\
1692 Tok Lane\\
3rd floor, Constitution Hall, Fairbanks, AK 99775\\
907-474-7300\\
\url{uaf-orca@alaska.edu}
\end{tabular}

\hfill  \scriptsize [syllabus version 1.01: \today] \normalsize

\end{document}

%%% Local Variables:
%%% mode: LaTeX
%%% TeX-master: t
%%% End:
